% rhythm: article float held here, standing between two paragraphs: paragraph to the float's picture, picture to its caption below, caption to the paragraph after; the caption's marker is the label the class generates, which opens its line
% boundary: par-float Figure above
% boundary: float-caption Figure below
% boundary: caption-par Cedar span
\documentclass{article}
\usepackage{tikz}
\usepackage{fontspec}
\setmainfont{OpenSans-Regular.ttf}[Path=../corpus/fonts/, BoldFont=OpenSans-Bold.ttf, ItalicFont=OpenSans-Italic.ttf, BoldItalicFont=OpenSans-BoldItalic.ttf]
\begin{document}
Alder is the paragraph above the float.

\begin{figure}[h]
\centering
\begin{tikzpicture}
\fill (0,0) rectangle (3,1.5);
\end{tikzpicture}
\caption{Birch caption}
\end{figure}

Cedar is the paragraph below the float.
\end{document}
